\subsection{OPERATIONS ON CARDINALS}

DEFINITION 3. Let \((\mathfrak{a}_t)_{t \in I}\) be a family of cardinals. The cardinal of the product (resp. sum) of the sets \(\mathfrak{a}_t\) is called the cardinal product (resp. cardinal sum) of the cardinals \(\mathfrak{a}_t\) and is denoted by \(\underset{i \in I}{\mathsf{P}} \mathfrak{a}_t\) ( \(\text{resp. } \sum_{i \in I} \mathfrak{a}_t\) ).

Whenever there is no risk of confusion we shall say simply ``product'' and ``sum'' in place of ``cardinal product'' and ``cardinal sum'', and we shall write \(\prod_{i\in I}a_{i}\) in place of \(P_{i\in I}a_{i}\) (cf.~Exercise 2).

PROPOSITION 4. Let \((\mathbf{E}_t)_{i \in I}\) be a family of sets, P their product, and S their sum, and let \(a_t\) be the cardinal of \(\mathbf{E}_t\) . Then the cardinal of P (resp. S) is the cardinal product (resp. cardinal sum) of the family \((a_t)_{i \in I}\) .

For there exists a bijection of P (resp. S) onto the product (resp. sum) of the sets \((\mathfrak{a}_{t})\) (Chapter II, § 4, no. 8, Proposition 10, and § 5, no. 7, Corollary to Proposition 11).

COROLLARY. If \((\mathbf{E}_{t})_{t\in\mathbf{I}}\) is any family of sets, the cardinal of the union \(\bigcup_{t\in I}E_{t}\) is at most equal to the sum \(\sum\) Card \((\mathbf{E}_{t})\) .

For there exists a mapping of the sum S of the \(E_{t}\) onto the union of the \(E_{t}\) (Chapter II, §4, no. 8); the Corollary therefore follows from Propositions 3 and 4.

PROPOSITION 5. (a) Let \((\mathfrak{a}_i)_{i\in I}\) be a family of cardinals, and let \(f\) be a bijection of a set \(\mathbf{K}\) onto the index set \(\mathbf{I}\) . Then

\[
\sum_ {x \in \mathbb {K}} a _ {f (x)} = \sum_ {i \in I} a _ {i}, \quad \underset {x \in \mathbb {K}} {\mathsf {P}} a _ {f (x)} = \underset {i \in I} {\mathsf {P}} a _ {i}.
\]

\begin{enumerate}
\def\labelenumi{(\alph{enumi})}
\setcounter{enumi}{1}
\tightlist
\item
  Let \((\mathfrak{a}_t)_{t\in \mathbf{I}}\) be a family of cardinals and let \((\mathrm{J}_{\lambda})_{\lambda \in \mathbf{L}}\) be a partition of \(\mathbf{I}\) . Then
\end{enumerate}

\[
\sum_ {i \in I} a _ {i} = \sum_ {\lambda \in L} \left(\sum_ {i \in J _ {\lambda}} a _ {i}\right), \quad P _ {i \in I} a _ {i} = P _ {\lambda \in L} \left(P _ {i \in J _ {\lambda}} a _ {i}\right)
\]

(``associativity of the sum and product'').

\begin{enumerate}
\def\labelenumi{(\alph{enumi})}
\setcounter{enumi}{2}
\tightlist
\item
  Let \(((\mathfrak{a}_{\lambda t})_{t\in J_{\lambda}}\lambda \in L\) be a family (indexed by L) of families of cardinals, and let \(I = \prod_{\lambda \in I}J_{\lambda}\) . Then
\end{enumerate}

\[
\underset {\lambda \in \mathbf {L}} {\mathbf {P}} \left(\sum_ {i \in \mathrm{J} _ {\lambda}} a _ {\lambda i}\right) \underset {f \in \mathbf {I}} {\sum} = \left(\underset {\lambda \in \mathbf {L}} {\mathbf {P}} a _ {\lambda , f (\lambda)}\right)
\]

(``distributivity of product over sum'').

The relations in (a) follow from the analogous formulae for the union and product of sets, for the fact that \(f\) is a bijection implies that if \((\mathbf{X}_t)_{i \in \mathbf{I}}\) is a family of mutually disjoint sets, the elements of the family \((\mathbf{X}_{f(\mathbf{x})})_{\mathbf{x} \in \mathbb{K}}\) are also mutually disjoint (cf.~Chapter II, §4, no. 1, Proposition 1 and §5, no. 3, Proposition 4).

¶ The relations in (b) are immediate consequences of the associativity formulae for unions and products (Chapter II, §4, no. 2, Proposition 2 and §5, no. 5, Proposition 7) and the distributivity of intersection over union (Chapter II, §5, no. 6, Proposition 8), which shows that if \((\mathbf{X}_{t})_{i\in I}\) is a family of mutually disjoint sets, then the elements of the family

\[
\left(\bigcup_ {i \in J _ {\lambda}} X _ {i}\right) _ {\lambda \in L}
\]

are also mutually disjoint.

Finally, (c) follows from the distributivity of the product over union and intersection (Chapter II, §5, no. 6, Proposition 9 and Corollary 1).

Let \(\mathfrak{a}\) and \(\mathfrak{b}\) be two cardinals. If \(\mathbf{I}\) is a set consisting of two distinct elements (e.g., the cardinal 2), there exists a mapping \(f\) of \(\mathbf{I}\) onto \(\{\mathfrak{a}, \mathfrak{b}\}\) which defines a family of cardinals. The sum and product of this family depend only on a and b (by reason of Proposition 5(a)); these cardinals are called respectively the sum and the product of a and b, and are denoted by \(\mathfrak{a} + \mathfrak{b}\) and \(\mathfrak{a}\mathfrak{b}\) . Similarly for the sum and product of three or more cardinals. Proposition 5 then implies the following corollary:

COROLLARY. Let a, b, c be cardinals. Then

\begin{enumerate}
\def\labelenumi{(\arabic{enumi})}
\tightlist
\item
\end{enumerate}

\[
\mathfrak {a} + \mathfrak {b} = \mathfrak {b} + \mathfrak {a}, \quad \mathfrak {a}\mathfrak {b} = \mathfrak {b}\mathfrak {a};\tag{2}
\]

\[
\mathfrak {a} + (\mathfrak {b} + \mathfrak {c}) = (\mathfrak {a} + \mathfrak {b}) + \mathfrak {c}, \quad \mathfrak {a}(\mathfrak {b}\mathfrak {c}) = (\mathfrak {a}\mathfrak {b})\mathfrak {c};\tag{3}
\]

\[
\mathfrak {a} (\mathfrak {b} + \mathfrak {c}) = \mathfrak {a}\mathfrak {b} + \mathfrak {a}\mathfrak {c}.
\]

